\documentclass[9pt,t]{beamer}
% =====================================================================
% finance-beamer.tex — shared Beamer style for the LLM-in-Finance course
% =====================================================================
% Reusable slide style. Each deck \input's this immediately after
% \documentclass{beamer}. It provides:
%   * the Madrid theme + book colour palette (matches the printed book)
%   * two book-style framing blocks:
%       \begin{bigpicture} ... \end{bigpicture}   ("the bigger picture")
%       \begin{underhood}   ... \end{underhood}    ("under the hood")
%     These mirror the book's `context` and `deepdive` tcolorboxes so the
%     same intuition-vs-mechanism scaffold the reader meets in the book
%     reappears on the slides.
%   * a Python listings style for code frames
%   * appendix conventions: \appendixstart drops a divider and switches
%     to non-numbered backup frames so "extra / less central" material
%     lives after the main talk without inflating the slide count.
%
% Usage:
%   \documentclass{beamer}
%   % =====================================================================
% finance-beamer.tex — shared Beamer style for the LLM-in-Finance course
% =====================================================================
% Reusable slide style. Each deck \input's this immediately after
% \documentclass{beamer}. It provides:
%   * the Madrid theme + book colour palette (matches the printed book)
%   * two book-style framing blocks:
%       \begin{bigpicture} ... \end{bigpicture}   ("the bigger picture")
%       \begin{underhood}   ... \end{underhood}    ("under the hood")
%     These mirror the book's `context` and `deepdive` tcolorboxes so the
%     same intuition-vs-mechanism scaffold the reader meets in the book
%     reappears on the slides.
%   * a Python listings style for code frames
%   * appendix conventions: \appendixstart drops a divider and switches
%     to non-numbered backup frames so "extra / less central" material
%     lives after the main talk without inflating the slide count.
%
% Usage:
%   \documentclass{beamer}
%   % =====================================================================
% finance-beamer.tex — shared Beamer style for the LLM-in-Finance course
% =====================================================================
% Reusable slide style. Each deck \input's this immediately after
% \documentclass{beamer}. It provides:
%   * the Madrid theme + book colour palette (matches the printed book)
%   * two book-style framing blocks:
%       \begin{bigpicture} ... \end{bigpicture}   ("the bigger picture")
%       \begin{underhood}   ... \end{underhood}    ("under the hood")
%     These mirror the book's `context` and `deepdive` tcolorboxes so the
%     same intuition-vs-mechanism scaffold the reader meets in the book
%     reappears on the slides.
%   * a Python listings style for code frames
%   * appendix conventions: \appendixstart drops a divider and switches
%     to non-numbered backup frames so "extra / less central" material
%     lives after the main talk without inflating the slide count.
%
% Usage:
%   \documentclass{beamer}
%   \input{../_style/finance-beamer.tex}
%   \title{...}\subtitle{...}\author{...}\institute{...}\date{\today}
%   \begin{document} ... \end{document}
% =====================================================================

\usetheme{Madrid}
\usecolortheme{whale}
\setbeamertemplate{navigation symbols}{}
\setbeamertemplate{caption}[numbered]
\setbeamercovered{transparent}

% --- Density controls: keep dense, equation-heavy frames on one slide ---
% Slightly smaller body text + tighter lists/blocks. Frame titles and the
% Madrid chrome keep their normal size, so the deck still reads as Madrid,
% just with more vertical room for content.
\setbeamerfont{normal text}{size=\small}
\setbeamertemplate{itemize/enumerate body begin}{\small}
\setbeamertemplate{itemize/enumerate subbody begin}{\footnotesize}
% Tighter block spacing.
\setbeamertemplate{block begin}{%
  \par\vskip2pt%
  \begin{beamercolorbox}[colsep*=.75ex,leftskip=1ex,rightskip=1ex]{block title}%
    \usebeamerfont*{block title}\insertblocktitle%
  \end{beamercolorbox}%
  {\parskip0pt\vskip-1pt}%
  \begin{beamercolorbox}[colsep*=.75ex,vmode,leftskip=1ex,rightskip=1ex]{block body}%
  \vskip1pt}
\setbeamertemplate{block end}{\end{beamercolorbox}\vskip2pt}

\usepackage{amsmath, amssymb}
\usepackage{booktabs}
\usepackage{xcolor}
\usepackage{listings}
\usepackage[skins, breakable]{tcolorbox}

% --- Book colour palette (kept in sync with book/preamble.tex) ---
\definecolor{linkblue}{RGB}{14, 75, 140}
\definecolor{deepdivebg}{RGB}{235, 244, 255}
\definecolor{narrativebg}{RGB}{255, 252, 240}
\definecolor{narrativerule}{RGB}{180, 130, 40}
\definecolor{steelblue}{RGB}{70, 130, 180}
\definecolor{forestgreen}{RGB}{34, 139, 34}
\definecolor{crimson}{RGB}{220, 20, 60}
\definecolor{darkorange}{RGB}{255, 140, 0}
\definecolor{codebg}{RGB}{248, 248, 248}
\definecolor{coderule}{RGB}{210, 210, 210}
\definecolor{keywordblue}{RGB}{0, 100, 180}
\definecolor{commentgray}{RGB}{120, 120, 120}
\definecolor{stringgreen}{RGB}{30, 130, 30}

% Tie the Madrid structural colour to the book's link blue.
\setbeamercolor{structure}{fg=linkblue}
\setbeamercolor{title}{fg=white}
\setbeamercolor{frametitle}{fg=white}

% --- "the bigger picture" : narrative / intuition framing -----------
\newtcolorbox{bigpicture}{%
  enhanced, breakable, boxsep=2pt,
  colback  = narrativebg,
  colframe = narrativerule,
  boxrule  = 0pt, toprule = 1.4pt, bottomrule = 0.4pt,
  leftrule = 0pt, rightrule = 0pt, arc = 0pt,
  left = 6pt, right = 6pt, top = 4pt, bottom = 5pt,
  fonttitle = \footnotesize\scshape\color{narrativerule},
  title = {the bigger picture}, attach title to upper = {\par\smallskip},
}

% --- "under the hood" : the mechanism / technical detail ------------
\newtcolorbox{underhood}{%
  enhanced, breakable, boxsep=2pt,
  colback  = deepdivebg,
  colframe = linkblue,
  boxrule  = 0pt, toprule = 1.4pt, bottomrule = 0.4pt,
  leftrule = 0pt, rightrule = 0pt, arc = 0pt,
  left = 6pt, right = 6pt, top = 4pt, bottom = 5pt,
  fonttitle = \footnotesize\scshape\color{linkblue},
  title = {under the hood}, attach title to upper = {\par\smallskip},
}

% --- Python code style ---------------------------------------------
\lstset{
  basicstyle    = \ttfamily\footnotesize,
  keywordstyle  = \color{keywordblue}\bfseries,
  commentstyle  = \color{commentgray}\itshape,
  stringstyle   = \color{stringgreen},
  showstringspaces = false,
  language      = Python,
  breaklines    = true,
  backgroundcolor = \color{codebg},
  frame         = single,
  rulecolor     = \color{coderule},
  numbers       = none,
  aboveskip     = 4pt, belowskip = 4pt,
}

% --- Appendix conventions ------------------------------------------
% \appendixstart{Title} : begins the "extra material" appendix. Frames
% after it are not counted toward the main deck's frame total, keeping
% the core talk lean while preserving the deeper / optional content.
\newcommand{\appendixstart}[1]{%
  \appendix
  \setbeamertemplate{frametitle continuation}{}%
  \begin{frame}[noframenumbering]
    \begin{center}
      {\Large\scshape\color{linkblue} Appendix}\\[4pt]
      {\large #1}\\[10pt]
      {\footnotesize Extra / optional material — beyond the core lecture.}
    \end{center}
  \end{frame}
}
% Use \begin{frame}[noframenumbering]{...} for each appendix frame.

%   \title{...}\subtitle{...}\author{...}\institute{...}\date{\today}
%   \begin{document} ... \end{document}
% =====================================================================

\usetheme{Madrid}
\usecolortheme{whale}
\setbeamertemplate{navigation symbols}{}
\setbeamertemplate{caption}[numbered]
\setbeamercovered{transparent}

% --- Density controls: keep dense, equation-heavy frames on one slide ---
% Slightly smaller body text + tighter lists/blocks. Frame titles and the
% Madrid chrome keep their normal size, so the deck still reads as Madrid,
% just with more vertical room for content.
\setbeamerfont{normal text}{size=\small}
\setbeamertemplate{itemize/enumerate body begin}{\small}
\setbeamertemplate{itemize/enumerate subbody begin}{\footnotesize}
% Tighter block spacing.
\setbeamertemplate{block begin}{%
  \par\vskip2pt%
  \begin{beamercolorbox}[colsep*=.75ex,leftskip=1ex,rightskip=1ex]{block title}%
    \usebeamerfont*{block title}\insertblocktitle%
  \end{beamercolorbox}%
  {\parskip0pt\vskip-1pt}%
  \begin{beamercolorbox}[colsep*=.75ex,vmode,leftskip=1ex,rightskip=1ex]{block body}%
  \vskip1pt}
\setbeamertemplate{block end}{\end{beamercolorbox}\vskip2pt}

\usepackage{amsmath, amssymb}
\usepackage{booktabs}
\usepackage{xcolor}
\usepackage{listings}
\usepackage[skins, breakable]{tcolorbox}

% --- Book colour palette (kept in sync with book/preamble.tex) ---
\definecolor{linkblue}{RGB}{14, 75, 140}
\definecolor{deepdivebg}{RGB}{235, 244, 255}
\definecolor{narrativebg}{RGB}{255, 252, 240}
\definecolor{narrativerule}{RGB}{180, 130, 40}
\definecolor{steelblue}{RGB}{70, 130, 180}
\definecolor{forestgreen}{RGB}{34, 139, 34}
\definecolor{crimson}{RGB}{220, 20, 60}
\definecolor{darkorange}{RGB}{255, 140, 0}
\definecolor{codebg}{RGB}{248, 248, 248}
\definecolor{coderule}{RGB}{210, 210, 210}
\definecolor{keywordblue}{RGB}{0, 100, 180}
\definecolor{commentgray}{RGB}{120, 120, 120}
\definecolor{stringgreen}{RGB}{30, 130, 30}

% Tie the Madrid structural colour to the book's link blue.
\setbeamercolor{structure}{fg=linkblue}
\setbeamercolor{title}{fg=white}
\setbeamercolor{frametitle}{fg=white}

% --- "the bigger picture" : narrative / intuition framing -----------
\newtcolorbox{bigpicture}{%
  enhanced, breakable, boxsep=2pt,
  colback  = narrativebg,
  colframe = narrativerule,
  boxrule  = 0pt, toprule = 1.4pt, bottomrule = 0.4pt,
  leftrule = 0pt, rightrule = 0pt, arc = 0pt,
  left = 6pt, right = 6pt, top = 4pt, bottom = 5pt,
  fonttitle = \footnotesize\scshape\color{narrativerule},
  title = {the bigger picture}, attach title to upper = {\par\smallskip},
}

% --- "under the hood" : the mechanism / technical detail ------------
\newtcolorbox{underhood}{%
  enhanced, breakable, boxsep=2pt,
  colback  = deepdivebg,
  colframe = linkblue,
  boxrule  = 0pt, toprule = 1.4pt, bottomrule = 0.4pt,
  leftrule = 0pt, rightrule = 0pt, arc = 0pt,
  left = 6pt, right = 6pt, top = 4pt, bottom = 5pt,
  fonttitle = \footnotesize\scshape\color{linkblue},
  title = {under the hood}, attach title to upper = {\par\smallskip},
}

% --- Python code style ---------------------------------------------
\lstset{
  basicstyle    = \ttfamily\footnotesize,
  keywordstyle  = \color{keywordblue}\bfseries,
  commentstyle  = \color{commentgray}\itshape,
  stringstyle   = \color{stringgreen},
  showstringspaces = false,
  language      = Python,
  breaklines    = true,
  backgroundcolor = \color{codebg},
  frame         = single,
  rulecolor     = \color{coderule},
  numbers       = none,
  aboveskip     = 4pt, belowskip = 4pt,
}

% --- Appendix conventions ------------------------------------------
% \appendixstart{Title} : begins the "extra material" appendix. Frames
% after it are not counted toward the main deck's frame total, keeping
% the core talk lean while preserving the deeper / optional content.
\newcommand{\appendixstart}[1]{%
  \appendix
  \setbeamertemplate{frametitle continuation}{}%
  \begin{frame}[noframenumbering]
    \begin{center}
      {\Large\scshape\color{linkblue} Appendix}\\[4pt]
      {\large #1}\\[10pt]
      {\footnotesize Extra / optional material — beyond the core lecture.}
    \end{center}
  \end{frame}
}
% Use \begin{frame}[noframenumbering]{...} for each appendix frame.

%   \title{...}\subtitle{...}\author{...}\institute{...}\date{\today}
%   \begin{document} ... \end{document}
% =====================================================================

\usetheme{Madrid}
\usecolortheme{whale}
\setbeamertemplate{navigation symbols}{}
\setbeamertemplate{caption}[numbered]
\setbeamercovered{transparent}

% --- Density controls: keep dense, equation-heavy frames on one slide ---
% Slightly smaller body text + tighter lists/blocks. Frame titles and the
% Madrid chrome keep their normal size, so the deck still reads as Madrid,
% just with more vertical room for content.
\setbeamerfont{normal text}{size=\small}
\setbeamertemplate{itemize/enumerate body begin}{\small}
\setbeamertemplate{itemize/enumerate subbody begin}{\footnotesize}
% Tighter block spacing.
\setbeamertemplate{block begin}{%
  \par\vskip2pt%
  \begin{beamercolorbox}[colsep*=.75ex,leftskip=1ex,rightskip=1ex]{block title}%
    \usebeamerfont*{block title}\insertblocktitle%
  \end{beamercolorbox}%
  {\parskip0pt\vskip-1pt}%
  \begin{beamercolorbox}[colsep*=.75ex,vmode,leftskip=1ex,rightskip=1ex]{block body}%
  \vskip1pt}
\setbeamertemplate{block end}{\end{beamercolorbox}\vskip2pt}

\usepackage{amsmath, amssymb}
\usepackage{booktabs}
\usepackage{xcolor}
\usepackage{listings}
\usepackage[skins, breakable]{tcolorbox}

% --- Book colour palette (kept in sync with book/preamble.tex) ---
\definecolor{linkblue}{RGB}{14, 75, 140}
\definecolor{deepdivebg}{RGB}{235, 244, 255}
\definecolor{narrativebg}{RGB}{255, 252, 240}
\definecolor{narrativerule}{RGB}{180, 130, 40}
\definecolor{steelblue}{RGB}{70, 130, 180}
\definecolor{forestgreen}{RGB}{34, 139, 34}
\definecolor{crimson}{RGB}{220, 20, 60}
\definecolor{darkorange}{RGB}{255, 140, 0}
\definecolor{codebg}{RGB}{248, 248, 248}
\definecolor{coderule}{RGB}{210, 210, 210}
\definecolor{keywordblue}{RGB}{0, 100, 180}
\definecolor{commentgray}{RGB}{120, 120, 120}
\definecolor{stringgreen}{RGB}{30, 130, 30}

% Tie the Madrid structural colour to the book's link blue.
\setbeamercolor{structure}{fg=linkblue}
\setbeamercolor{title}{fg=white}
\setbeamercolor{frametitle}{fg=white}

% --- "the bigger picture" : narrative / intuition framing -----------
\newtcolorbox{bigpicture}{%
  enhanced, breakable, boxsep=2pt,
  colback  = narrativebg,
  colframe = narrativerule,
  boxrule  = 0pt, toprule = 1.4pt, bottomrule = 0.4pt,
  leftrule = 0pt, rightrule = 0pt, arc = 0pt,
  left = 6pt, right = 6pt, top = 4pt, bottom = 5pt,
  fonttitle = \footnotesize\scshape\color{narrativerule},
  title = {the bigger picture}, attach title to upper = {\par\smallskip},
}

% --- "under the hood" : the mechanism / technical detail ------------
\newtcolorbox{underhood}{%
  enhanced, breakable, boxsep=2pt,
  colback  = deepdivebg,
  colframe = linkblue,
  boxrule  = 0pt, toprule = 1.4pt, bottomrule = 0.4pt,
  leftrule = 0pt, rightrule = 0pt, arc = 0pt,
  left = 6pt, right = 6pt, top = 4pt, bottom = 5pt,
  fonttitle = \footnotesize\scshape\color{linkblue},
  title = {under the hood}, attach title to upper = {\par\smallskip},
}

% --- Python code style ---------------------------------------------
\lstset{
  basicstyle    = \ttfamily\footnotesize,
  keywordstyle  = \color{keywordblue}\bfseries,
  commentstyle  = \color{commentgray}\itshape,
  stringstyle   = \color{stringgreen},
  showstringspaces = false,
  language      = Python,
  breaklines    = true,
  backgroundcolor = \color{codebg},
  frame         = single,
  rulecolor     = \color{coderule},
  numbers       = none,
  aboveskip     = 4pt, belowskip = 4pt,
}

% --- Appendix conventions ------------------------------------------
% \appendixstart{Title} : begins the "extra material" appendix. Frames
% after it are not counted toward the main deck's frame total, keeping
% the core talk lean while preserving the deeper / optional content.
\newcommand{\appendixstart}[1]{%
  \appendix
  \setbeamertemplate{frametitle continuation}{}%
  \begin{frame}[noframenumbering]
    \begin{center}
      {\Large\scshape\color{linkblue} Appendix}\\[4pt]
      {\large #1}\\[10pt]
      {\footnotesize Extra / optional material — beyond the core lecture.}
    \end{center}
  \end{frame}
}
% Use \begin{frame}[noframenumbering]{...} for each appendix frame.


\title{Portfolio Optimization \& Quantitative Trading with LLMs}
\subtitle{Lecture 10 --- Signals, Backtests, and Risk}
\author{Juan F.\ Imbet}
\institute{Paris Dauphine -- PSL University}
\date{\today}

\begin{document}

\begin{frame}\titlepage\end{frame}

\begin{frame}{Outline}\tableofcontents\end{frame}

% ============================================================
\section{Classical Portfolio Theory and Its Limits}
% ============================================================

\begin{frame}{The Central Question --- and a Stubborn Paradox}
\begin{bigpicture}
Given $N$ risky assets, how should a rational investor allocate wealth?
Markowitz (1952) reduced this to a quadratic program: the \emph{efficient frontier}.
Decades of machinery followed --- yet a naive $1/N$ equal-weight portfolio
\textbf{routinely beats} the ``optimal'' one (DeMiguel et al., 2009).
\end{bigpicture}

\medskip
\textbf{Why does more sophistication underperform?}
\begin{itemize}
  \item The theory is right; the \emph{inputs} are noisy.
  \item Estimating $\boldsymbol{\mu}$ and $\boldsymbol{\Sigma}$ from data dominates the gains.
  \item This is the opening for LLMs: extract \emph{more precise} $\boldsymbol{\mu}$ from text.
\end{itemize}

\medskip
\begin{block}{Thesis of the lecture}
LLMs are \textbf{information-extraction engines}, not optimizers. They produce
signals; classical theory dictates how those signals become portfolio weights.
\end{block}
\end{frame}

\begin{frame}{Markowitz Mean-Variance Optimization}
\begin{underhood}
Let $\mathbf{r}\in\mathbb{R}^N$ be excess returns, $\boldsymbol{\mu}=\mathbb{E}[\mathbf{r}]$,
$\boldsymbol{\Sigma}=\mathrm{Cov}[\mathbf{r}]$ (positive definite). For weights $\mathbf{w}$:
\[
\mu_p = \boldsymbol{\mu}^\top\mathbf{w}, \qquad \sigma_p^2 = \mathbf{w}^\top\boldsymbol{\Sigma}\,\mathbf{w}.
\]
\end{underhood}

\medskip
\textbf{Efficient portfolio} at target return $\bar\mu$:
\[
\min_{\mathbf{w}} \; \mathbf{w}^\top\boldsymbol{\Sigma}\,\mathbf{w}
\quad\text{s.t.}\quad
\boldsymbol{\mu}^\top\mathbf{w}=\bar\mu,\;\; \mathbf{1}^\top\mathbf{w}=1,\;\; \mathbf{w}\ge\mathbf{0}.
\]

\textbf{Maximum-Sharpe (tangency) portfolio}, unconstrained closed form:
\[
\mathbf{w}^{\text{SR}} \;\propto\; \boldsymbol{\Sigma}^{-1}\boldsymbol{\mu}.
\]
Scales \emph{inversely} with risk, \emph{linearly} with the return signal.
\end{frame}

\begin{frame}{Why Estimation Error Wins}
\textbf{The inputs are the problem:}
\begin{itemize}
  \item Sample $\hat{\boldsymbol{\mu}}$ is noisy: with 60 months of data the SE of a
        monthly mean is $\approx\hat\sigma/\sqrt{60}$ --- on the order of the mean itself.
  \item Sample $\hat{\boldsymbol{\Sigma}}$ is near-singular when $N$ is large vs.\ $T$;
        its inverse \emph{amplifies} errors.
  \item Result: heavy bets on recently lucky assets $\to$ unstable, poor out-of-sample.
\end{itemize}

\medskip
\textbf{Standard remedies:}
\begin{itemize}
  \item \emph{Shrinkage} --- pull $\hat{\boldsymbol{\Sigma}}$ toward a structured target.
  \item \emph{Robust optimization} --- replace point estimates with uncertainty sets.
  \item \emph{Factor models}: $\boldsymbol{\Sigma}=\mathbf{B}\mathbf{F}\mathbf{B}^\top+\mathbf{D}$,
        reducing $N(N{+}1)/2$ entries to $K$ factors (Fama--French 3- / 5-factor).
\end{itemize}

\medskip
\begin{alertblock}{The DeMiguel et al.\ (2009) finding}
Across \textbf{14 datasets}, $1/N$ beats sophisticated optimizers on out-of-sample
Sharpe. Precise $\boldsymbol{\mu}$ is the missing ingredient --- enter LLMs.
\end{alertblock}
\end{frame}

\begin{frame}{Black--Litterman: Equilibrium Prior $+$ Views}
\begin{bigpicture}
Don't trust raw sample means. Start from a prior that reflects \emph{market equilibrium},
then update it with investor views via Bayes. This decouples \emph{where} you have a
view from \emph{what} the view is --- exactly the shape an LLM signal takes.
\end{bigpicture}

\begin{underhood}
Equilibrium implied returns from market-cap weights $\mathbf{w}^{\text{mkt}}$, risk aversion $\delta$:
\[
\boldsymbol{\Pi}=\delta\,\boldsymbol{\Sigma}\,\mathbf{w}^{\text{mkt}},
\qquad \boldsymbol{\mu}\sim\mathcal{N}(\boldsymbol{\Pi},\tau\boldsymbol{\Sigma}).
\]
$K$ views $\mathbf{P}\boldsymbol{\mu}=\mathbf{q}+\boldsymbol{\varepsilon}$,
$\boldsymbol{\varepsilon}\sim\mathcal{N}(\mathbf{0},\boldsymbol{\Omega})$. Posterior mean:
\[
\hat{\boldsymbol{\mu}}=
\bigl[(\tau\boldsymbol{\Sigma})^{-1}+\mathbf{P}^\top\boldsymbol{\Omega}^{-1}\mathbf{P}\bigr]^{-1}
\bigl[(\tau\boldsymbol{\Sigma})^{-1}\boldsymbol{\Pi}+\mathbf{P}^\top\boldsymbol{\Omega}^{-1}\mathbf{q}\bigr].
\]
\end{underhood}
No views $\Rightarrow \hat{\boldsymbol{\mu}}\to\boldsymbol{\Pi}$, optimum reduces to the market portfolio (sanity check).
\end{frame}

\begin{frame}{Black--Litterman: A Single LLM View}
\begin{exampleblock}{Example --- Apple absolute view}
An LLM reads Apple's earnings call $\to$ $0.72$ probability of an EPS beat.
Translate to a BL view:
\[
\mathbf{P}=\mathbf{e}_{\text{AAPL}}^\top,\quad q=1\%\ \text{(monthly excess)},\quad
\Omega_{11}=(0.5\%)^2.
\]
Prior $\Pi_{\text{AAPL}}=0.3\%$. The posterior shifts upward by an amount
$\propto (q-\mathbf{P}\boldsymbol{\Pi})$ relative to total uncertainty
$\tau\mathbf{P}\boldsymbol{\Sigma}\mathbf{P}^\top+\boldsymbol{\Omega}$.
\end{exampleblock}

\begin{itemize}
  \item High confidence ($\Omega_{11}$ small) $\Rightarrow$ posterior $\approx q=1\%$.
  \item Low confidence $\Rightarrow$ posterior barely leaves $\Pi_{\text{AAPL}}$.
  \item BL prevents \emph{any single noisy signal} from dominating the portfolio.
\end{itemize}

\medskip
\footnotesize Empirical validation: Mantshimuli \& Mwamba (2025) aggregate multi-LLM
sentiment views via an LSTM into BL --- improving Sharpe and annualised return over
single-LLM and market-cap benchmarks. Yuksel (2025) treats the LLM as an
evolutionary portfolio \emph{selector}.
\end{frame}

\begin{frame}{Where LLMs Can --- and Cannot --- Add Value}
\begin{columns}[T]
\begin{column}{0.5\textwidth}
\begin{block}{\textcolor{forestgreen}{Clear value}}
\begin{itemize}
  \item Structuring \emph{alternative data} (calls, filings) into factors/views at scale
  \item \emph{Risk monitoring} --- flagging emerging risks before they hit the statements
  \item \emph{Analyst-view aggregation} with calibrated uncertainty
\end{itemize}
\end{block}
\end{column}
\begin{column}{0.5\textwidth}
\begin{alertblock}{No clear value}
\begin{itemize}
  \item Return prediction beyond a few days (factor risk dominates)
  \item Replacing factor models (not built to estimate $\boldsymbol{\Sigma}$)
  \item Timing in isolation (no order-flow / microstructure access)
\end{itemize}
\end{alertblock}
\end{column}
\end{columns}

\medskip
\textbf{Two hypotheses underpin the case:}
\begin{itemize}
  \item \emph{Information}: text holds un-impounded return information
        (Grossman--Stiglitz: efficiency cannot be perfect). Evidence \emph{mixed} ---
        Lopez-Lira \& Tang (2023) find significant but small, costly-to-harvest signal.
  \item \emph{Extraction}: LLMs read negation, hedging, sarcasm better than bag-of-words
        --- this conjunct is largely \emph{unambiguous} (FinBERT).
\end{itemize}
\end{frame}

% ============================================================
\section{LLMs as Alternative Data Processors}
% ============================================================

\begin{frame}{Text: The Least-Saturated Alternative Data}
\begin{bigpicture}
Satellite imagery and card flows erode fast once adopted. \textbf{Text} ---
earnings calls, analyst reports, filings, news, central-bank speak --- generates
terabytes of natural language a year and stays comparatively un-mined.
LLMs convert this prose into numerical signals at scale.
\end{bigpicture}

\medskip
\textbf{Two construction pathways in this section:}
\begin{enumerate}
  \item \textbf{Views} for Black--Litterman --- a per-firm $q_i$ with calibrated uncertainty.
  \item \textbf{Factors} for Fama--French --- a long-short ``language factor'' tested for alpha.
\end{enumerate}
\end{frame}

\begin{frame}{Extracting Analyst Views from Text}
\textbf{Four-stage extraction pipeline:}
\begin{enumerate}
  \item \textbf{Retrieve \& chunk} --- document from feed / SEC EDGAR, split semantically.
  \item \textbf{Score} --- LLM (or FinBERT) assigns sentiment per chunk.
  \item \textbf{Aggregate} --- position-weighted into a document signal $s_i\in[-1,1]$.
  \item \textbf{Normalize} --- within industry \& period to remove sector trends.
\end{enumerate}

\begin{underhood}
Map signal to a BL view (uncertainty \emph{inverse} to signal strength):
\[
q_i=\alpha\cdot s_i+\Pi_i,\qquad \Omega_{ii}=\sigma_q^2\cdot(1-|s_i|).
\]
$\alpha$ calibrated by regressing past signals on realized excess returns.
Strong signal ($|s_i|\!\approx\!1$) $\to$ low uncertainty; neutral $\to$ prior dominates.
\end{underhood}

\medskip
\textbf{Beware staleness / look-ahead:} use only text available \emph{before} the
rebalance close. After-hours summaries leaking into next-day trades $=$ look-ahead bias.
\end{frame}

\begin{frame}{Worked: Earnings-Call Views (100 S\&P 500 names)}
\begin{exampleblock}{Prompt}
\textit{``You are a sell-side equity analyst. Given this earnings-call excerpt, assign
a return sentiment score from $-1$ (bearish) to $+1$ (bullish) over the next month.
Output only the number.''}
\end{exampleblock}

\begin{itemize}
  \item CFO ``raising full-year guidance 5\%, no demand softening'' $\to$ score $\in[0.6,0.9]$.
  \item CEO ``cautious on macro, right-sizing cost structure'' $\to$ score $\in[-0.6,-0.3]$.
\end{itemize}

\medskip
\textbf{Calibrated results} ($T=800$):
\begin{itemize}
  \item Signal vs.\ one-month excess return correlation $\approx 0.08$ (significant),
        implying an information ratio $\approx 0.25$.
  \item Inserted into a BL optimizer ($\delta=3$, $\tau=0.05$): Sharpe
        \textbf{15\% above $1/N$} --- \emph{before} transaction costs.
\end{itemize}

\medskip
\footnotesize Meskovskis \& Kenyon (2024): an LLM-generated \emph{SWOT} view from strategic
filings supports longer horizons where price-based backtesting lacks power.
\end{frame}

\begin{frame}{Constructing LLM-Based Factors}
\begin{underhood}
Factor decomposition of returns:
\[
r_{i,t}=\alpha_i+\sum_{k=1}^{K}\beta_{ik}f_{k,t}+\varepsilon_{i,t}.
\]
Fama--French: MKT, SMB, HML ($+$ RMW, CMA in the 5-factor model).
\end{underhood}

\medskip
\textbf{Language factor (LF).} Demean $s_{i,t}$ cross-sectionally, sort into quintiles:
\[
\text{LF}_t = R_t^{\text{top quintile}} - R_t^{\text{bottom quintile}}.
\]

\textbf{Does it earn alpha beyond the classics?} Test $\hat a>0$ in
\[
\text{LF}_t=a+\beta_{\text{MKT}}\text{MKT}_t+\beta_{\text{SMB}}\text{SMB}_t
+\beta_{\text{HML}}\text{HML}_t+\beta_{\text{RMW}}\text{RMW}_t+\beta_{\text{CMA}}\text{CMA}_t+\varepsilon_t.
\]
Earnings-call alpha exists \emph{short-run} ($\le 5$ days) then decays as the market reads the same text.
\end{frame}

\begin{frame}{The Factor Zoo --- and Automated Factor Discovery}
\begin{alertblock}{Multiple-testing inflation (Harvey, Liu \& Zhu, 2016)}
With hundreds of candidate factors tested on one dataset, $t>2$ produces \emph{many}
false discoveries. A factor surviving \textbf{Bonferroni} or a multiple-comparison
Sharpe test is far stronger evidence than raw $p<0.05$.
\end{alertblock}

\medskip
\textbf{LLMs can automate the search itself:}
\begin{itemize}
  \item \emph{AlphaQuant} (Yuksel, 2025): an LLM-evolutionary loop that proposes,
        evaluates, and refines alpha features --- replacing manual feature engineering.
  \item This amplifies the multiple-testing problem: the backtesting discipline that
        follows becomes \textbf{non-negotiable}.
\end{itemize}
\end{frame}

% ============================================================
\section{LLM-Guided Algorithmic Trading}
% ============================================================

\begin{frame}{Where in the Trading Stack Do LLMs Belong?}
\begin{bigpicture}
Algorithmic trading spans HFT (microseconds, microstructure) to systematic macro
(weeks--months, fundamentals). LLMs matter in the \textbf{slow} regime: where signal
\emph{quality}, not delivery \emph{speed}, decides performance. Latency in hours, not ms.
\end{bigpicture}

\medskip
A signal $s_t$ drives a position $w_t=f(s_t)$ that must:
\begin{itemize}
  \item scale with signal strength,
  \item be bounded (no runaway concentration),
  \item decay as the signal ages.
\end{itemize}
\end{frame}

\begin{frame}{Sentiment-Driven Signal Generators}
\begin{underhood}
Z-score the raw LLM output over a rolling window (e.g.\ 252 days), then size:
\[
z_{i,t}=\frac{\tilde s_{i,t}-\hat\mu_s}{\hat\sigma_s},
\qquad
w_{i,t}=\frac{z_{i,t}}{\sum_j|z_{j,t}|}\cdot W_{\max}.
\]
Balanced $\pm$ signals $\Rightarrow$ dollar-neutral long-short (insulated from the market).
\end{underhood}

\medskip
\textbf{Sentiment decays} --- model it explicitly (exponential / EWMA):
\[
s_{i,t}=\sum_{\tau\le t}\tilde s_{i,\tau}\,e^{-\lambda(t-\tau)},
\quad \lambda\approx 0.1\text{--}0.5\ \text{day}^{-1}\ \text{(half-life 1--7 days)}.
\]

\medskip
\footnotesize Kirtac \& Germano (2024): GPT-3-class news scores predict next-day returns
with above-chance directional accuracy and a high in-sample post-cost Sharpe ---
but single-sample Sharpes rarely replicate out of sample. A \emph{magnitude} baseline,
not a promise.
\end{frame}

\begin{frame}{LLM-Guided RL for Order Execution}
\textbf{Almgren--Chriss optimal execution} (linear impact):
\[
C(\mathbf{x})=\tfrac{1}{2}\gamma X^2+\eta\sum_{k}\tau_k v_k^2+\lambda\sigma^2\sum_{k}\tau_k x_k^2.
\]
\footnotesize permanent impact $+$ temporary slippage $+$ shortfall-variance penalty;
front-load when risk-averse, spread uniformly when not.

\normalsize
\medskip
\textbf{RL extends this} (FinRL): state $s_t=$ (inventory, time left, spread, momentum);
action $a_t\in[0,X_t]$; reward $r_t=-\text{cost}(a_t)$.

\begin{underhood}
LLM role $=$ \textbf{state augmentation}, not replacement:
\[
\tilde s_t=[\,s_t^{\text{micro}},\ \text{LLM-alert}_t\,],\qquad
\text{LLM-alert}_t\in\{-1,0,+1\}.
\]
A negative-surprise alert triggers \emph{accelerated} execution, cutting adverse selection.
\end{underhood}
\footnotesize HARLF (Coriat \& Benhamou, 2025): three-tier hierarchical RL injecting LLM
sentiment into the \emph{reward} --- strong risk-adjusted returns on 2018--2024 data.
\end{frame}

\begin{frame}{Why LLMs Are \emph{Not} Stock Pickers}
\begin{alertblock}{Statistically significant $\ne$ economically profitable}
The market-efficiency literature has drawn this line for decades.
\end{alertblock}

\begin{itemize}
  \item \textbf{Semi-strong EMH} (Fama, 1970): public info is already in prices.
        Large caps adjust to news in \emph{minutes} --- faster than an LLM API can read it.
  \item \textbf{Grossman--Stiglitz}: \emph{some} return must accrue to informed trading,
        but if everyone subscribes to the same API the signal \emph{crowds} and alpha compresses.
  \item \textbf{Calibration problem}: a ``70\% positive'' LLM probability is \emph{not}
        a claim the stock rises 70\% of the time --- requires OOS calibration to returns.
\end{itemize}

\medskip
\begin{block}{The right niche}
\textbf{Heterogeneous, low-coverage, novel} sources --- small-cap filings,
emerging-market news, niche disclosures --- where analyst coverage is thin and
price discovery is slow. There the Grossman--Stiglitz return survives.
\end{block}
\end{frame}

% ============================================================
\section{Backtesting with LLM Signals}
% ============================================================

\begin{frame}{Backtesting: Necessary, Not Sufficient}
\begin{bigpicture}
A backtest asks: would this strategy have made money historically? A \emph{realistic}
one accounts for frictions, data limits, and the multiple-comparison problem.
Single-period in-sample Sharpe is nearly useless --- the model has seen the future.
\end{bigpicture}

\medskip
\textbf{Walk-forward validation} = time-series cross-validation. For fold $m$ with step $\Delta$:
\begin{itemize}
  \item \emph{Train}: $[m\Delta,\ m\Delta+T_{\text{train}})$
  \item \emph{Validate}: next $T_{\text{val}}$ (hyperparameters only)
  \item \emph{Test}: next $\Delta$ --- report performance \textbf{only} on concatenated test windows.
\end{itemize}
For LLM signals the ``parameters'' include $\alpha$, decay $\lambda$, scale $W_{\max}$ ---
\emph{never} tuned on the test period.
\end{frame}

\begin{frame}{Out-of-Sample Sharpe and Transaction Costs}
\textbf{Primary metric --- annualised OOS Sharpe:}
\[
\text{SR}^{\text{OOS}}=\frac{\bar r_p^{\text{OOS}}}{\hat\sigma_p^{\text{OOS}}}\cdot\sqrt{252}.
\]
\footnotesize Rule of thumb: $\text{SR}>1.5$ over $\ge 3$ years OOS before paper trading.

\normalsize
\begin{underhood}
Turnover cost model, $\Delta w_{i,t}=w_{i,t}-w_{i,t-1}$:\quad
$\text{TC}_t=c\sum_i|\Delta w_{i,t}|$. Daily strategy, 15\% one-way turnover, $c=0.05\%$:
\[
0.15\times 252\times 0.05\%\approx \textbf{1.89\%}\ \text{p.a.} \ \Rightarrow\ \text{erases marginal alpha.}
\]
\end{underhood}

\textbf{Cost components:} bid--ask (single-digit bps large-cap, tens small-cap),
market impact ($\eta x^2+\gamma x$), financing ($>5\%$ p.a.\ hard-to-borrow).

\textbf{Capacity:} trade $\le 5$--$10\%$ of ADV; earnings-call signals crowd in the first hours.
\end{frame}

\begin{frame}{Look-Ahead and Survivorship: The Killers}
\small
\textbf{Look-ahead bias --- subtle LLM-specific forms:}
\begin{enumerate}
  \item \textbf{Embedding leakage} --- model fine-tuned on transcripts \emph{post-dating} the trade.
  \item \textbf{Model-selection leakage} --- best LLM variant chosen by full-period performance.
  \item \textbf{Timestamp errors} --- calls release after close; day-$t$ open uses unavailable info.
\end{enumerate}

\textbf{Survivorship bias:} today's S\&P 500 list excludes the distressed names removed ---
exactly where a risk strategy drew down. Use \emph{point-in-time} constituents.

\begin{underhood}
\textbf{Backtest overfitting} --- PBO (Bailey et al., 2014) $\to 1$ for null strategies as more
combinations are tried. Minimum track record length (Lopez de Prado, 2018):
\[
T^*=1+\Bigl(1-\gamma_1\text{SR}^*+\tfrac{\gamma_2-1}{4}(\text{SR}^*)^2\Bigr)\frac{Z_{1-\alpha}^2}{(\text{SR}^*)^2}.
\]
Fat-tailed LLM-sentiment returns $\Rightarrow$ $T^*$ of \emph{several years}.
\end{underhood}
\end{frame}

% ============================================================
\section{Risk Management Applications}
% ============================================================

\begin{frame}{Tail Risk from Textual Signals}
\begin{bigpicture}
Qualitative warnings --- management caution, litigation, sovereign risk --- shift the
\emph{tail} of the return distribution \textbf{before} they show up in prices. LLMs read
this stream and feed it into volatility and tail-loss estimates in near real time.
\end{bigpicture}

\textbf{CVaR (Expected Shortfall)} at level $\alpha$, a coherent, LP-minimisable
risk measure (Rockafellar \& Uryasev, 2000):
\[
\text{CVaR}_\alpha(R)=-\mathbb{E}\bigl[R\mid R\le q_\alpha(R)\bigr].
\]

\begin{underhood}
Textual volatility signal $v_t$ enters a GARCH(1,1):\quad
$\sigma_t^2=\omega+\varphi_1\varepsilon_{t-1}^2+\beta_1\sigma_{t-1}^2+\kappa\,v_{t-1}$.
$\kappa>0$: elevated risk-language forecasts higher variance --- fatter tails next period.
Connects to NVIX (Manela \& Moreira, 2017), now at the \emph{firm} level.
\end{underhood}
\end{frame}

\begin{frame}{CVaR-Constrained Optimization \& Live Monitoring}
\textbf{De-risk automatically when the text turns dark:}
\[
\min_{\mathbf{w}}\ \mathbf{w}^\top\boldsymbol{\Sigma}_t\mathbf{w}
\quad\text{s.t.}\quad \text{CVaR}_\alpha(\mathbf{w}^\top\mathbf{r}_t)\le\bar C,
\]
$\boldsymbol{\Sigma}_t$ time-varying (DCC-GARCH). High $v_{t-1}$ in a sector $\Rightarrow$
$\sigma_{i,t}^2\uparrow\Rightarrow$ tighter CVaR $\Rightarrow$ automatic de-risking.

\medskip
\textbf{Monitoring pipeline --- three layers:}
\begin{itemize}
  \item \textbf{Ingestion} --- poll EDGAR, news APIs, social at configurable frequency.
  \item \textbf{Classification} --- LLM emits structured JSON
        (\texttt{risk\_level}, \texttt{risk\_type}, \texttt{affected\_positions}).
  \item \textbf{Action} --- route critical alerts to a risk officer or a pre-approved auto-response.
\end{itemize}

\medskip
\footnotesize Calibration: tune to the firm's portfolio; beat \emph{alert fatigue}
($<20$/day for 100 positions); match latency to use (minutes EOD, seconds intraday).
\end{frame}

\begin{frame}[fragile]{Example: An 8-K Alert That Saved Half the Loss}
\textbf{4:17 PM} --- a held name files an 8-K announcing an SEC investigation into
revenue recognition. Within 30s the LLM classifies:
\begin{verbatim}
{ "risk_level": "critical", "risk_type": "legal",
  "affected_positions": ["TICKER_X"],
  "summary": "SEC formal order into revenue recognition.
              Historical precedent: 20-40% downside risk." }
\end{verbatim}
Action layer cuts the position $2.1\%\to1.0\%$ at the next open.
\textbf{Next morning the stock opens $-28\%$.}

\medskip
\begin{alertblock}{Promise and limit}
The system acted faster than any human. But it could not foresee the exact
magnitude, nor whether the probe ends in no action or a restatement.
It \emph{reduces} tail exposure --- human judgment still sets the thresholds.
\end{alertblock}
\end{frame}

% ============================================================
\section{Summary}
% ============================================================

\begin{frame}{Lecture 10 --- Five Takeaways}
\small
\begin{enumerate}
  \item \textbf{Optimizers, not replaced.} Markowitz and Black--Litterman remain the
        right machinery; LLMs sharpen $\boldsymbol{\mu}$ from text. Extractors, not optimizers.
  \item \textbf{Factor construction is the most defensible channel.} Test LLM-derived
        long-short alpha vs.\ Fama--French; correct for multiple comparisons.
  \item \textbf{Efficiency limits but doesn't eliminate alpha.} Grossman--Stiglitz $\Rightarrow$
        value lives in low-coverage markets; large caps are the hardest target.
  \item \textbf{Backtesting discipline is non-negotiable.} Walk-forward, real costs,
        point-in-time data, overfitting correction --- research artefact vs.\ deployable strategy.
  \item \textbf{Real-time risk monitoring} is high-value, low-competition: structured risk
        signals $+$ GARCH augmentation $+$ continuous surveillance, with auditable value.
\end{enumerate}

\medskip
\begin{block}{Next}
Practical 10: build the efficient frontier, map an LLM score to a BL view, and run a
cost-aware walk-forward backtest in Python.
\end{block}
\end{frame}

% ============================================================
\appendixstart{Extended Derivations, Benchmarks, and Further Reading}
% ============================================================

\begin{frame}[noframenumbering]{Appendix: Sharpe-Ratio Geometry of the Frontier}
The unconstrained tangency portfolio $\mathbf{w}^{\text{SR}}\propto\boldsymbol{\Sigma}^{-1}\boldsymbol{\mu}$
maximises
\[
\frac{\boldsymbol{\mu}^\top\mathbf{w}}{\sqrt{\mathbf{w}^\top\boldsymbol{\Sigma}\mathbf{w}}}.
\]
\begin{itemize}
  \item The frontier (\texttt{gen\_efficient\_frontier.py}, deterministic, no data/network)
        places the tangency and $1/N$ portfolios \emph{almost on top of one another} ---
        a concrete picture of the DeMiguel et al.\ result.
  \item Adding the non-negativity constraint $\mathbf{w}\ge\mathbf{0}$ removes the closed form;
        the mean-variance program must be solved numerically as a QP.
\end{itemize}
\end{frame}

\begin{frame}[noframenumbering]{Appendix: Recent LLM-Portfolio Literature Map}
\small
\begin{tabular}{lll}
\toprule
\textbf{Work} & \textbf{Mechanism} & \textbf{Result} \\
\midrule
Mantshimuli \& Mwamba (2025) & Multi-LLM $\to$ LSTM $\to$ BL views & Sharpe / return $\uparrow$ vs.\ single-LLM \\
Yuksel (2025), \emph{alpha}  & LLM evolutionary tournament       & Portfolio selection, large universes \\
Yuksel (2025), AlphaQuant    & LLM-evolutionary factor search    & Automated alpha-feature discovery \\
Coriat \& Benhamou (2025)    & HARLF: 3-tier hierarchical RL     & Strong risk-adj.\ returns 2018--2024 \\
Meskovskis \& Kenyon (2024)  & SWOT views from strategic filings & Long-horizon view generation \\
Saha, Lyu et al.\ (2025)     & Survey of agentic PM systems      & Taxonomy $+$ open benchmarks \\
\bottomrule
\end{tabular}

\medskip
\normalsize Practitioner survey: Mann (2024) --- durable value sits in the
\emph{information-extraction layer}, not the decision layer.
\end{frame}

\begin{frame}[noframenumbering]{Appendix: News, Volatility, and the EMH Evidence}
\begin{itemize}
  \item \textbf{Tetlock (2007)}: dictionary sentiment from a newspaper column predicted
        excess returns over a decade ago --- the \emph{incremental} value of LLMs over
        dictionaries is the relevant comparison; broad-coverage signals have decayed.
  \item \textbf{NVIX} (Manela \& Moreira, 2017): text-based disaster-risk measure from
        newspaper archives predicted market variance and crash risk. LLMs push this to
        firm-level, real-time.
  \item \textbf{Regulatory context}: Fed SR 11-7 requires model outputs be explainable,
        validated, auditable --- directly binding on LLM risk signals (see the RegTech chapter).
\end{itemize}
\end{frame}

\begin{frame}[noframenumbering]{Appendix: Signal-to-View Calibration Detail}
\textbf{Why $\Omega_{ii}=\sigma_q^2(1-|s_i|)$?}
\begin{itemize}
  \item It makes view \emph{uncertainty} shrink as signal \emph{magnitude} grows.
  \item Neutral signals ($s_i\approx 0$) get high $\Omega_{ii}$, so the equilibrium prior
        $\Pi_i$ dominates --- the optimizer ignores noise.
  \item $\alpha$ is the \emph{economic translation} from a unitless score to a return view,
        estimated by regressing realized excess returns on lagged signals.
\end{itemize}
\medskip
\textbf{Pitfall:} calibrating $\alpha$ on the test window is model-selection leakage ---
it must use only data predating the evaluation window.
\end{frame}

\end{document}
